\clearpage
\VerdierPhysicalPageAnchor{0160}%
\begingroup\small\setlength{\parskip}{0pt}\setlength{\emergencystretch}{1.5em}
\VerdierSemanticAnchor{ch3-3-2-7-condition}
\noindent\textit{such that the diagram}
\[
X_1\xrightarrow{u_1}X_2\xrightarrow{f_3u_2}Y_3
\]
\textit{can be included in a distinguished triangle
$(X_1,X_2,Y_3,u_1,f_3u_2,u_3)$.}
This follows immediately from the proof of Proposition (3.2.5).

\VerdierSemanticAnchor{ch3-3-2-8}
\medskip\noindent\textbf{3.2.8. Proof of assertions a) and b) of Theorem
3.2.1.}\quad
We have already shown (3.2.6) that $\mathrm A(\mathcal D)$ is abelian. Let
$F\colon\mathcal D\to\mathcal D'$ be an exact functor. It follows
immediately from Proposition (3.2.7) that $\mathrm A(F)$ is left exact.
Moreover, it is clear from the construction of $\mathrm A(\mathcal D)$ that
the category $\mathrm A(\mathcal D^\circ)^\circ$ is canonically isomorphic
to $\mathrm A(\mathcal D)$ and that the functor
$\mathrm A(F^\circ)^\circ$ is canonically isomorphic to $\mathrm A(F)$.
Since $F^\circ$ is exact (1.1.7), it follows that
$\mathrm A(F^\circ)$ is left exact and consequently that $\mathrm A(F)$ is
also right exact, hence exact.

The admissible objects of $\mathcal D^\wedge$ are isomorphic to subobjects
and quotient objects of representable functors. Consequently, by (3.2.3),
every object of $\mathrm A(\mathcal D)$ is isomorphic to a subobject and to
a quotient object of an object of $\mathcal D$. Representable functors are
projective objects in $\mathcal D^\wedge$. Therefore, by (3.2.5), they are
projective in the full subcategory of admissible objects of
$\mathcal D^\wedge$. Consequently, by (3.2.3), the objects of $\mathcal D$
are projective in $\mathrm A(\mathcal D)$. It follows that the objects of
$\mathcal D^\circ$ are projective in $\mathrm A(\mathcal D^\circ)$ and
hence, on passing to opposite categories, that the objects of $\mathcal D$
are injective in $\mathrm A(\mathcal D)$. The formal proof that
$\mathrm A_{\mathcal D,\mathcal D'}$ is fully faithful may safely be left to
the reader: a morphism between exact functors is uniquely determined once
its values on sufficiently many projective objects are known.

\VerdierSemanticAnchor{ch3-3-2-9}
\medskip\noindent\textbf{3.2.9. Completion of the proof of Theorem 3.2.1.}
We first show that the functor
$\mathcal H\colon\mathcal D\to\mathrm A(\mathcal D)$ is a cohomological
functor. Let $(X,Y,Z,f,g,h)$ be a distinguished triangle in $\mathcal D$.
We must prove that the sequence
\[
\mathcal H(X)\xrightarrow{\mathcal H(f)}\mathcal H(Y)
\xrightarrow{\mathcal H(g)}\mathcal H(Z).
\]
is exact. For this ((3.2.3) and (3.2.5)), it suffices to show that the sequence
\VerdierSemanticAnchor{ch3-3-2-9-1}
\[
h(X)\xrightarrow{h(f)}h(Y)\xrightarrow{h(g)}h(Z)
\tag{3.2.9.1}
\]
\endgroup
